\chapter{The link to the C code}\label{chap:ccode}

Lean has no semantics of C, and no theorem of this blueprint speaks about the C code of ratpoints:
the theorems speak about the model, which mirrors the C code component by component, with the same
loops, tables, bit arrays of 64-bit words and index arithmetic. This chapter describes what ties
the model to the C code of release~3.1.1: the record of which declaration models which lines of C,
and its checks; the lines that no declaration models; the departures of the model from the C code;
the differential tests of the program \verb|ratpoints-lean| against the C program; the builds with
other values of the compile-time constants; and CI. The nodes of the chapters before this one cite
the main ranges of their definitions; Appendix~\ref{chap:correspondence}, generated from the
record, lists every range with its owners and their deviations, or its exclusion.

The link is not a proof. That a definition does what its lines of C do is checked by reading the
two side by side; the record makes sure that every line of C is accounted for, that each claim is
made for a definite text on both sides, and that a change on either side is noticed; the
differential tests check the correspondence on behaviour, on finitely many inputs. Comparator
certifies the statements about the model and says nothing about the C code
(Section~\ref{sec:comparator}).

\section{The record}\label{sec:ccodeRecord}

Every line number refers to one commit of ratpoints, the pin: release~3.1.1, tag \verb|v3.1.1|,
commit \verb|1a99853|, named in \verb|correspondence/ratpoints.toml| and, for the links of the
blueprint, in the macro \verb|\ratpointscommit|. The record covers \verb|find_points.c|,
\verb|sift.c|, \verb|init.c|, \verb|sturm.c|, the headers \verb|rp-private.h|, \verb|ratpoints.h|
and \verb|primes.h|, the generators of tables \verb|gen_find_points_h.c| and
\verb|gen_init_sieve_h.c|, and the program \verb|main.c|. Every line of code of these files, that
is, every line that holds something besides white space and comments (a line of the preprocessor
counts), lies in exactly one recorded range, and no two ranges overlap.

A range is modelled when the docstring of a public declaration of the model names it, in a
paragraph that begins with \verb|C:|; an optional paragraph that begins with \verb|Deviation:| says
where the declaration departs from the C code. The docstring of \lname{Ratpoints.Sieve.squaresMod},
for instance, ends with
\begingroup\small
\begin{verbatim}
C: `find_points.c:37`, `gen_find_points_h.c:69-88`

Deviation: a function of `p`, computed when it is called, for the row of `p` in the table
`squares` that the build generates and that `find_points.c` includes. -/
\end{verbatim}
\endgroup
A range may have several owners (line~37 of \verb|find_points.c| includes three generated tables
and is a range of the three functions that compute them), and a declaration that calls another,
which models a part of its C code, cites the lines around that part. A private declaration carries
no record, since no node can name it. A range is excluded when a table \verb|[[excluded]]| of the
record gives it, with a kind and the reason (Section~\ref{sec:ccodeExclusions}).

The script \verb|scripts/correspondence.py| checks the record. It reads the docstrings of the model
with the help of \verb|scripts/Correspondence.lean|, and the C files at the pin from a clone of
ratpoints, and it checks:
\begin{enumerate}
  \item that the pin is a commit of the clone, to which its tag points where the clone has the tag,
    that every range lies within the lines of a file of the record, that no line lies in two ranges,
    and that every exclusion has a known kind and a reason;
  \item that every range and every C link of the blueprint is locked: for every range, the lock
    holds the SHA-256 hash of its text at the pin and its owners (the declarations that model it, or
    its exclusion), and for every link the hash of the text that it shows; a range whose text or
    owners differ from those of the lock, and a new link, are reported until they have been compared
    with their counterparts by hand and locked again by \verb|scripts/correspondence.py lock|, which
    writes the lock into \verb|correspondence/lock.json|;
  \item that the appendix is the one that the record generates, and that every link inside a node
    of the dependency graph that names declarations with \verb|\lean| is one of the ranges of the
    model declarations that the node names, exactly as recorded (a node need not cite all of
    them; prose and remarks may link to any range);
  \item with \verb|--complete|, which CI passes, that every line of code lies in a range.
\end{enumerate}
After a move of the pin, a range counts as locked when some locked range has its text and its
owners, and a link when some locked link shows its text, wherever their lines are; so a move of the
pin to another release of ratpoints makes the check report exactly the ranges whose text or owners
changed and the links whose text changed (\verb|correspondence/README.md| describes the procedure).

\section{The exclusions}\label{sec:ccodeExclusions}

An excluded range has one of the kinds \verb|threads|, \verb|memory|, \verb|instrumentation|,
\verb|interface| and \verb|other|.

\verb|threads|: the pool of threads, with which ratpoints sieves the blocks of its loop over the
denominators apart (\csource{find_points.c}{4129}{4519}), its set-up in the search, the atomic loads
and stores of the fields that the threads share (\csource{rp-private.h}{76}{84}), the fields of the
arguments that hold the number of threads and the pool, and the default of one thread in
\verb|main.c|. The model sieves the blocks in turn. The thread lemma (Theorem~\ref{thm:threads})
shows that blocks sifted apart and delivered in their order give the same; that the pool of the C
code meets its hypotheses is read off the C code (Chapter~\ref{chap:theorems}).

\verb|memory|: the buffers of the sieve entries, the rows, the tables, the words of the forbidden
divisors and the prime powers, with their alignment and enlargement
(\csource{find_points.c}{181}{270}) and their release (\csource{find_points.c}{377}{416}), the
arrays of the run and of a worker, the copy of the caller's intervals in the Sturm step, the GMP
variables that \verb|main.c| initializes and releases, and the failures for lack of memory,
error~10 of the program among them. The model allocates nothing itself and cannot fail for lack of
memory; such failures inside modelled ranges, such as the drop of the forbidden divisors beyond
the table whose words find no memory, are deviations of their owners.

\verb|instrumentation|: the output under \verb|DEBUG|, the timers and counters of
\verb|RP_PHASE_TIMING| and \verb|RP_PHASE_COUNTS| (\csource{sift.c}{55}{176}), the report of
\verb|RP_PRIME_STATS|, \verb|RP_STOP_AFTER|, \verb|RP_MOD_CHOICE|, \verb|RP_MOD_COUNTS|, the check
\verb|RP_VERIFY_MODULI| of the rows (\csource{init.c}{299}{353}) and two probes for a benchmark,
all compiled only on request or empty in the default build. Output under \verb|DEBUG| inside a
modelled block goes with the block, whose deviation says that the model prints nothing there; the
report under \verb|RATPOINTS_VERBOSE| (\verb|-v|) is modelled (Chapter~\ref{chap:output}).

\verb|interface|: what only another caller of the library meets: \verb|find_points_init|
(\csource{find_points.c}{272}{375}) and its prototype, the input fields of the arguments that
\verb|find_points_work| saves and puts back, the checks of arrays that only such a caller can fail
(\csource{find_points.c}{4618}{4623}), the declaration of the arguments in \verb|main.c| and the
initialization of the fields that the model passes as values, and the prototype of the Sturm step.

\verb|other|: lines without meaning for the model: the includes, those of the vector intrinsics
among them, and the prototypes; the inlining attributes and the pragma that unrolls the loops of
\verb|tail_leg| (\csource{sift.c}{466}{472}); \verb|FLOOR|, \verb|RATPOINTS_DEFAULT_SP1| and
\verb|RATPOINTS_DEFAULT_SP2|, which the C code reads nowhere, and a count stored that nothing reads;
the text of message~1, which \verb|main| never prints (\csource{main.c}{977}{978}); the frames of
the \verb|main| functions of the generators of tables; and the arm of \verb|exact_scale| for a build
without a 128-bit integer type (\csource{find_points.c}{3239}{3256}), since the default build
compiles the arm with \verb|unsigned __int128|.

\section{The deviations}\label{sec:ccodeDeviations}

The \verb|Deviation:| paragraphs record, declaration by declaration, where the model departs from
its lines of C. Most of the departures are of the following kinds. None corrects the C code: the
model is that of release~3.1.1 as it is.

\emph{Integer types.} The model computes with natural numbers and integers where the C code has
\verb|long|s, \verb|int|s, \verb|unsigned long|s and GMP integers, and with rational numbers for GMP
rationals; where a value of the C code could leave its type, the range obligation is a lemma that
the deviation names. Lean's 64-bit words stand where the C code depends on the words themselves: in
the bit arrays, the reductions by a reciprocal and the counters of the run, which wrap as
\verb|unsigned long|s wrap; in the test for common factors (\lname{Ratpoints.Sieve.relprime}), whose
step without branches acts on the two's complement of a negative \verb|long| and shifts it
arithmetically, as gcc defines these operations, which the C standard leaves to the implementation;
and in the bits of the numerators of an interval, the forms of the ends of the intervals and the
numerators of the Sturm step, whose arithmetic wraps where that of the \verb|long|s would overflow,
on no path that the search takes. The 128-bit products are taken on two words, the high word from
four products of 32-bit halves (\lname{Ratpoints.Sieve.mulHi}). The height bound is a natural
number, refused from~$2^{63}$ on with error~4, as the program refuses such a height when it reads
it.

\emph{Pointers and buffers.} Arrays and indices stand for pointers into arrays, an option for a
pointer that may be null, and values returned, often structures, for what the C code writes
through pointers or into buffers of its caller; state that the C code reaches through pointers,
such as the delivery of the points, the tables of rows (\lname{Ratpoints.Sieve.RowStore}) and the
counters, is passed and returned. A search gives back what the program reads
(Definition~\ref{def:Options}), not the other fields that the C code leaves in the arguments. The
end of a row that is not built yet is~$p$ (\lname{Ratpoints.Sieve.setRows}), where \verb|set_rows|
adds~$p$ to the null pointer (\csource{find_points.c}{3149}{3170}), which the C standard leaves
undefined; this happens in ordinary runs, whenever \verb|set_rows| reports a row as missing, and the
end is set again when the row is built, before it is read.

\emph{Generated tables.} The squares, the inverses and the rows of~$0$ of the primes, which the
build generates into \verb|find_points.h|, are functions of the prime, computed when they are
called: \lname{Ratpoints.Sieve.squaresMod}, \lname{Ratpoints.Sieve.inversesMod} and
\lname{Ratpoints.Sieve.sieves0}. The functions that the build generates into \verb|init_sieve.h|,
one for each prime, which build its rows, are two functions, \lname{Ratpoints.Sieve.sieveInit1} and
\lname{Ratpoints.Sieve.sieveInit2}, with the prime an argument; and a pointer to such a function is
the kind of the builder, with a match on it (Chapter~\ref{chap:rows}).

\emph{Vector code.} A bit array of the default build is a vector of $256$~bits; in the model it is
$\mathrm{rbaPack} = 4$ consecutive words of an array of \verb|UInt64|, and the \verb|&| of two bit
arrays the \verb|&| of their words (\lname{Ratpoints.Sieve.andBitArray}). The registers of the
chunked first stage are a vector of $\mathrm{chunk} = 16$ bit arrays, the code that the build
unrolls under \verb|#if (RATPOINTS_CHUNK >= k)| is a loop over them
(\lname{Ratpoints.Sieve.chunkSieve}), and the conditional compilation that chooses an arm is a test
on the constant (Chapter~\ref{chap:sieve}). The other builds differ in these constants
(Section~\ref{sec:ccodeConstants}).

\emph{Doubles.} The doubles of the search region are their IEEE~754 bits (\lname{Ratpoints.Double}),
with a function, which owns no range of C, for each operation that ratpoints applies to them, as gcc
and glibc compute it with rounding to nearest; \verb|strtod| and \verb|printf| under \verb|%f| are
modelled by their values, not by glibc's algorithms (Chapter~\ref{chap:region}). The doubles of the
cost model are Lean's \verb|Float|, with the platform's arithmetic, \verb|log| and \verb|pow|, and
no proof reads their values; the sorts by \verb|qsort| are a stable merge sort
(\lname{Ratpoints.sortFrom}), which gives the order of glibc's \verb|qsort|, a stable merge sort
whenever it can allocate its buffer, since no key is a NaN (Chapter~\ref{chap:moduli}). Where the C
code computes in \verb|double| what the model computes exactly (the comparison of $2pH$ with the
limit of the third stage, the integer square roots of the shapes), the deviation shows that the
outcome is the same.

\emph{Structure and the program.} The model has one function, \lname{Ratpoints.findPointsVerbose},
for \verb|find_points|, \verb|find_points_work| and the start of \verb|find_points_work_1|, whose
rest is \lname{Ratpoints.findPointsWork}; parts of long functions and bodies of loops are functions
of their own; arrays that the C code fills once per denominator under a flag are thunks; and the
report under \verb|RATPOINTS_VERBOSE| goes through a print channel of the search. The program writes
its output once, at the end, where \verb|main.c| prints after every point; it differs from ratpoints
on a closed pipe and on format words that are not UTF-8, where the runtime of Lean acts otherwise
than that of C (Chapter~\ref{chap:output}).

\emph{Inputs that ratpoints never meets.} Where the C code would overflow or divide by zero on an
input that it is never given, the model takes a defined value, and the contract excludes the input
or the proofs show that it does not occur: the negation of~$-2^{63}$ in the test for common factors
wraps, and the integer square root rounded up gives~$1$ at~$0$. The loop that counts the trailing
zeros does not end on~$0$, as in the C code (\lname{Ratpoints.Sieve.ctz}), and the masks of the
forbidden valuations drop the bound $m < 63$ of the C code, which never binds for a height that is a
\verb|long| (Chapter~\ref{chap:denominators}).

\section{The differential tests}\label{sec:ccodeTests}

The script \verb|scripts/difftest.py| runs \verb|ratpoints-lean|, the model compiled by Lean
(Chapter~\ref{chap:output}), on the test suites of ratpoints and on further command lines, and
compares it with the outputs that the suites record and with the C program at the pin on the same
command line; the C program must be the default build, without the machine-dependent tuning that
\verb|make tune| writes. Every list of points is compared in its order. Under \verb|-q| the points
are compared with a recorded output and byte for byte with the output of the C program, and so are
the outputs with the options and formats of the output and the reports, with messages 3, 4 and~8 on
the search region, the moduli used and the search intervals, the banner exchanged for that of the C
program, also under \verb|-v|, where the report of the library comes between the points. The
outputs that \verb|test4.sh| filters are compared after the filter and as a whole; those that
\verb|test3.sh| filters, as a whole, before the filter. Under \verb|-x|, where the program prints
the pairs that survive the sieve, the pairs are compared byte for byte with the C program, since the
model chooses and corrects its moduli as ratpoints does, and, without a format and without
\verb|-v|, they also pass a check that does not depend on the moduli: normalized coprime pairs in
the range of the height and in the search region, none repeated, and the points among them the
same, in the same order, in both outputs. Exit statuses and messages are compared exactly, but for
the banner and the usage; \verb|ratpoints-lean| must write nothing to its standard error (a panic
of the Lean runtime ends it); and failed writes are compared too.

The default suites search up to the height~$30000$, and above it in narrow search regions. Since the
choice of the moduli computes with the \verb|log| and \verb|pow| of the platform's C library, which
glibc does not round correctly and chooses by the processor, the comparisons that depend on it (the
pairs under \verb|-x|, message~4 and the report under \verb|-v|) are made with both programs on one
machine; CI runs both on its runner. On the state of the project that this blueprint describes, the
suites make these comparisons, none of which fails:
\begin{center}
\small
\begin{tabular}{@{}lp{0.68\textwidth}r@{}}
  suite & what it runs & comparisons \\
  \hline
  \verb|test4| & the runs of \verb|test4.sh|, against \verb|testbase4| and the C program, but
    for~$13$ above the height~$30000$ & $590$ \\
  \verb|test3| & the runs of \verb|test3.sh| & $21$ \\
  \verb|rptest| & the curves of \verb|testdata.h|, \verb|testdata-many.h| and
    \verb|testdata-degrees.h| at the height~$100$, against their recorded points & $1206$ \\
  \verb|random| & random polynomials of degree~$1$ to~$8$ at heights up to~$30$ & $300$ \\
  \verb|outputs| & random options of the output, formats and strings & $300$ \\
  \verb|errors| & hand-written command lines for every check of \verb|main.c| and every option,
    the refusals, the report under \verb|-v|, searches stopped at once at the height~$10^{12}$,
    failed writes & $384$ \\
  \verb|fuzz| & random command lines & $800$ \\
  \verb|params| & the options that set the moduli and the cut points & $272$ \\
  \verb|high| & point-rich curves up to the height~$30000$, and curves on which the C program
    corrects its moduli & $19$ \\
  \verb|regions| & random search regions and ranges of the denominators & $400$ \\
  \verb|depths| & random depths of the Sturm step & $300$ \\
  \verb|tuning| & random values of the tuning options & $300$ \\
  \verb|threads| & the C program with \verb|-t 4| and \verb|-t 7| & $205$ \\
  \verb|verbose| & random searches under \verb|-v| & $300$ \\
  \hline
  & the default suites & $5397$ \\
  \hline
  \verb|test4high| & the runs of \verb|test4.sh| at the heights $40000$ to~$100000$ & $13$ \\
  \verb|adapt| & searches at the heights $16383$ to~$50000$ on which the C program corrects its
    moduli, under \verb|-x| with and without \verb|-A| & $92$ \\
\end{tabular}
\end{center}
The last two suites are not run by default, nor the comparison of a build of the C program with a
fixed length of the blocks (\verb|-DRP_BLOCK_LEN=3|) with \verb|ratpoints-lean| built with the same
length (\lname{Ratpoints.blockLen}), which tests the boundaries of the blocks. The \verb|#guard|
examples of the model, which every build runs, compare components with values taken from the C
program, gcc and glibc, none of which depends on the \verb|log| or \verb|pow| of the platform.

\section{Other values of the constants}\label{sec:ccodeConstants}

The file \verb|Ratpoints/Model/Constants.lean| holds the constants that the C code fixes when it is
compiled, each a definition with the value of the default build and a docstring that names its
macro: among them the bound on the degree (\lname{Ratpoints.maxDegree}), the number of bits of the
moduli, the width of a bit array (\lname{Ratpoints.rbaPack}), the chunk of the first stage
(\lname{Ratpoints.chunk}), the default depth of the Sturm step, the limits of the reductions by a
reciprocal, the doubles that tune the cost model and the constants of the blocks
(\lname{Ratpoints.adaptLag}, \lname{Ratpoints.blockLen}). No proof depends on their values beyond a
few facts that the file decides: that the width, the chunk and the array sizes are positive, that
the limits of the reductions are at most~$2^{32}$, that a bit array has $2^{\mathrm{rbaShift}}$ bits
with $\mathrm{rbaShift} < 32$, that $\mathrm{maxBitsInPrime} \le 10$, that the bound on the degree
is positive and below~$2^{57}$, that a fixed length of the blocks is positive, and a few more of
that kind.

The script \verb|scripts/constants_check.py| checks this by building copies of the committed
project, with all its proofs and without the \verb|#guard| examples of the model, which test the
default values, for eleven other sets of values, each of which builds: two that change every constant, one to small values (among them bit
arrays of~$64$~bits, the chunk~$1$, the lag~$0$ and blocks of length~$1$) and one to large values
(bit arrays of~$512$~bits, the chunk~$17$, which compiles the other arm of the first stage, the
lag~$1000$); four that change the chunk alone, to~$2$, $3$, $5$ and~$12$; three that pair the
chunked first stage with the widths of the other builds, $512$, $128$ and~$64$~bits; one with the
lag~$1$ and blocks of length~$3$; and one with one block wanted. So the proofs hold for the models
of the configurations of the C code that differ in these constants, the widths of the bit arrays of
all its builds among them. The script builds the copies one after the other, with the built
packages of the repository, and runs on demand, outside CI; the sets of values are a table in it,
and \verb|--variant| gives others.

\section{Continuous integration}\label{sec:ccodeCI}

The workflow \verb|.github/workflows/build-project.yml| runs on every push to the main branch and on
every pull request. It checks the layout (\verb|scripts/check_layout.py|: the imports of each part;
no \verb|partial|, \verb|unsafe|, \verb|implemented_by|, \verb|extern| or experimental option in the
model, which would let the executable run code that the proofs do not see; and nowhere a tactic that
adds an axiom, such as \verb|native_decide|), the style of the Lean code and the length of its
lines, and the listing \verb|Ratpoints.lean| of the modules; builds, with every warning a failure,
the model, the specification and the proofs, \verb|ratpoints-lean| and the solution of comparator,
and then the challenge, whose \verb|sorry|s warn by design; runs the tests of the scripts; checks
the record with \verb|--complete| against ratpoints at the pin; and builds the C program at the pin
and runs the default suites of the differential tests against it. The build fails if a main theorem
depends on an axiom beyond \verb|propext|, \verb|Classical.choice| and \verb|Quot.sound|, by the
guards \verb|#print axioms| under \verb|#guard_msgs|. Comparator (Section~\ref{sec:comparator}),
the suites \verb|test4high| and \verb|adapt|, the pair with a fixed length of the blocks and the
builds with other values of the constants (\verb|scripts/constants_check.py|) run outside CI.
